\section{Conclusion}

The suspension keyed to student number LXXTRI004 was identified as

\begin{equation}
    G(s) = \frac{1}{s^2 + \gap{b}\,s + \gap{k}},
\end{equation}

a \gap{underdamped / overdamped} second order system with gain \gap{A} m/V, damping ratio
\gap{zeta} and natural frequency \gap{wn} rad/s, giving a spring coefficient of \gap{k} and
a damping coefficient of \gap{b}.

The step test supplied all three parameters and the frequency response test confirmed the
natural frequency independently, to within \gap{difference}. \gap{Say in one sentence
whether the two experiments agree, and if they do not, which one you trust and why.}

The main limitation is the 100~ms sample period, which is coarse compared with a damped
period of \gap{Td} s and can hide the true peak of the response. \gap{Add any limitation
your own data actually showed, for example a non-dominant pole that never became visible, or
a steady state that had not fully settled inside the record.}
